%% ============================================================
%%  SECCIÓN 7 — Resultados Experimentales
%% ============================================================

\section{Resultados Experimentales}
\label{sec:resultados}

En esta sección se presentan los datos obtenidos durante las mediciones
en el laboratorio. Para cada material se reportan primero los parámetros
geométricos iniciales y, a continuación, las mediciones registradas bajo
la aplicación de cargas progresivas. En el caso de la liga elástica se
incluyen también las mediciones correspondientes a la fase de descarga.

%% ------------------------------------------------------------------
\subsection{Resorte Metálico}
%% ------------------------------------------------------------------

Se determinaron los parámetros geométricos iniciales del resorte metálico
y su masa en reposo. La Tabla~\ref{tab:param_resorte} presenta dichos
parámetros.

\begin{table}[H]
    \centering
    \caption{Parámetros iniciales del resorte metálico.}
    \label{tab:param_resorte}
    \begin{tabular}{lcc}
        \toprule
        \textbf{Magnitud} & \textbf{Símbolo} & \textbf{Valor} \\
        \midrule
        Masa del resorte       & $m_r$ & $46{,}3  \pm 0{,}5\  \text{g}$  \\
        Longitud natural       & $L_0$ & $16{,}60 \pm 0{,}05\ \text{cm}$ \\
        Diámetro de la sección & $D_0$ & $14{,}00 \pm 0{,}02\ \text{mm}$ \\
        \bottomrule
    \end{tabular}
\end{table}

Las mediciones de longitud $L$ y diámetro $D$ del resorte bajo carga
progresiva se registran en la Tabla~\ref{tab:datos_resorte}. Se emplearon
seis configuraciones de masa, con valores nominales entre \qty{500}{g}
y \qty{1750}{g}. Para cada configuración se consignaron la masa nominal,
la masa real medida con la balanza, la longitud total del resorte y el
diámetro de su sección transversal deformada, necesario para el cálculo
del esfuerzo real $\sigma = F/S$.

\begin{table}[H]
    \centering
    \caption{Mediciones de longitud y diámetro del resorte bajo carga progresiva.
             Incertidumbres: $\Delta L = \pm 0{,}05\ \text{cm}$,\;
             $\Delta D = \pm 0{,}02\ \text{mm}$.}
    \label{tab:datos_resorte}
    \begin{tabular}{crrcc}
        \toprule
        \multirow{2}{*}{\textbf{N$^{\circ}$}}
            & \multicolumn{2}{c}{\textbf{Masa $m$ (g)}}
            & \textbf{$L$ (cm)}
            & \textbf{$D$ (mm)} \\
        \cmidrule(lr){2-3}
            & Nominal & Real & & \\
        \midrule
        1 &  500 &  499{,}8 & 29{,}8 & 14{,}22 \\
        2 &  750 &  754{,}3 & 32{,}1 & 14{,}14 \\
        3 & 1000 & 1005{,}0 & 35{,}0 & 14{,}04 \\
        4 & 1250 & 1259{,}5 & 38{,}5 & 13{,}98 \\
        5 & 1500 & 1504{,}8 & 41{,}7 & 13{,}94 \\
        6 & 1750 & 1759{,}3 & 44{,}6 & 13{,}90 \\
        \bottomrule
    \end{tabular}
\end{table}

%% ------------------------------------------------------------------
\subsection{Liga Elástica}
%% ------------------------------------------------------------------

Se determinaron los parámetros geométricos iniciales de la liga elástica
y su masa en reposo. La Tabla~\ref{tab:param_liga} presenta dichos
parámetros.

\begin{table}[H]
    \centering
    \caption{Parámetros iniciales de la liga elástica.}
    \label{tab:param_liga}
    \begin{tabular}{lcc}
        \toprule
        \textbf{Magnitud} & \textbf{Símbolo} & \textbf{Valor} \\
        \midrule
        Masa de la liga       & $m_l$ & $46{,}3  \pm 0{,}5\  \text{g}$   \\
        Longitud natural      & $L_0$ & $34{,}20 \pm 0{,}05\ \text{cm}$  \\
        Ancho de la sección   & $a_0$ & $14{,}80 \pm 0{,}02\ \text{mm}$  \\
        Espesor de la sección & $b_0$ & $5{,}50  \pm 0{,}02\ \text{mm}$  \\
        \bottomrule
    \end{tabular}
\end{table}

Las mediciones de longitud $L$, ancho $a$ y espesor $b$ de la sección
transversal se registran en la Tabla~\ref{tab:datos_liga}. Las mediciones
se realizaron en dos etapas: carga y descarga, registrándose las
dimensiones de la sección transversal deformada en cada caso, necesarias
para el cálculo del esfuerzo real $\sigma = F/S$. Se emplearon seis
configuraciones de masa en la fase de carga (entre \qty{500}{g} y
\qty{1750}{g}) y cinco en la fase de descarga (hasta \qty{500}{g}),
consignando en cada punto la masa nominal, la masa real, la longitud
total y las dimensiones de la sección transversal.

\begin{table}[H]
    \centering
    \caption{Mediciones de longitud y sección transversal de la liga en
             fases de carga y descarga.
             Incertidumbres: $\Delta L = \pm 0{,}25\ \text{cm}$,\;
             $\Delta a = \Delta b = \pm 0{,}02\ \text{mm}$.}
    \label{tab:datos_liga}
    \begin{tabular}{lcrrccc}
        \toprule
        \multirow{2}{*}{\textbf{Fase}}
            & \multirow{2}{*}{\textbf{N$^{\circ}$}}
            & \multicolumn{2}{c}{\textbf{Masa $m$ (g)}}
            & \textbf{$L$ (cm)}
            & \textbf{$a$ (mm)}
            & \textbf{$b$ (mm)} \\
        \cmidrule(lr){3-4}
            & & Nominal & Real & & & \\
        \midrule
        \multirow{6}{*}{\textbf{Carga}}
            & 1 &  500 &  499{,}8 & 36{,}7 & 14{,}24 & 5{,}60 \\
            & 2 &  750 &  754{,}3 & 38{,}6 & 13{,}76 & 5{,}46 \\
            & 3 & 1000 & 1005{,}0 & 41{,}5 & 13{,}14 & 5{,}24 \\
            & 4 & 1250 & 1259{,}5 & 45{,}0 & 13{,}14 & 5{,}24 \\
            & 5 & 1500 & 1504{,}8 & 47{,}1 & 13{,}12 & 5{,}14 \\
            & 6 & 1750 & 1759{,}3 & 48{,}3 & 12{,}84 & 5{,}14 \\
        \midrule
        \multirow{5}{*}{\textbf{Descarga}}
            & 5 & 1500 & 1504{,}8 & 47{,}4 & 12{,}32 & 5{,}12 \\
            & 4 & 1250 & 1259{,}5 & 46{,}0 & 13{,}30 & 5{,}40 \\
            & 3 & 1000 & 1005{,}0 & 45{,}5 & 13{,}42 & 5{,}42 \\
            & 2 &  750 &  754{,}3 & 42{,}2 & 13{,}52 & 5{,}36 \\
            & 1 &  500 &  499{,}8 & 39{,}4 & 13{,}60 & 5{,}58 \\
        \bottomrule
    \end{tabular}
\end{table}